% Source: https://artofproblemsolving.com/community/c3h2040636_a_question
% Stable identity: aops-topic-2040636-q01
\textbf{False.} Consider the system of \(n=2\) equations in \(2\) unknowns
\[
x+y=1,\qquad 2x+2y=2 .
\]
The second equation is twice the first, so every pair \((x,y)=(t,1-t)\), \(t\) real, is a solution: there are infinitely many solutions.

In general a system of \(n\) linear equations in \(n\) unknowns has exactly one solution if its coefficient determinant is nonzero; if the determinant is zero it has either no solution (e.g. \(x+y=1\), \(x+y=2\)) or infinitely many solutions (as above). So the possible numbers of solutions are \(0\), \(1\), or infinitely many.
